\section{The Definition Crisis}
\label{sec:definitions}

AI law has no shortage of definitions. It has a shortage of distinctions among the jobs those definitions are asked to perform. A legislature defines a regulated class to allocate compliance obligations. An engineer defines a system boundary to conduct a hazard analysis. A company describes a product to attract customers and limit contractual exposure. A court classifies an object to apply an inherited rule. Each act of definition selects features for a purpose. Trouble follows when the answer produced for one purpose is carried into another without argument.

The central example is personification. ``Agent'' is an entirely serviceable term for software that selects and sequences operations toward an objective. It is also a term of art in the law of representation, where an agent acts on behalf of a principal and can alter legal relations. ``Companion'' can describe a product's intended use. It can also evoke a reciprocal social relationship. ``Assistant'' can identify an interface while implying judgment, loyalty, or expertise. Nothing requires engineers and marketers to abandon these useful words. Courts must simply decline their hidden premise: a functional metaphor does not determine the system's legal position.

\subsection{A Field Organized Around Four Questions}

An emerging field becomes coherent when its recurring problems can be stated together, not when every dispute receives the same rule. Cyberlaw developed by asking how network architecture altered familiar institutions of control, attribution, privacy, property, and speech.\lrfootnote{See generally \artcite{lessigLawHorse1999}; \artcite{caloRobotics2015}; \artcite{balkinRobotics2015}.} AI law is reaching a similar point. The systems now appear in employment, education, welfare administration, medicine, finance, creative production, intimate communication, military procurement, and ordinary consumer transactions. The governing doctrines differ, but four questions recur:

\begin{table}[htbp]
\centering
\caption{The Four Questions of AI Regulation}
\label{tab:four-questions}
\small
\begin{tabularx}{\textwidth}{>{\raggedright\arraybackslash}p{0.16\textwidth} >{\raggedright\arraybackslash}p{0.31\textwidth} X}
\toprule
Question & Legal function & Recurring AI issue \\
\midrule
Who is regulated? & Identifies legal subjects and duty holders. & Provider, deployer, professional, public body, user, or component supplier. \\
What is regulated? & Identifies the object, activity, transaction, or risk. & Model, application, output, decision process, data practice, or represented service. \\
How is it regulated? & Specifies conduct rules and institutional controls. & Prohibition, disclosure, testing, human review, documentation, licensing, or procurement condition. \\
What follows from breach? & Connects violation to remedy and responsibility. & Damages, injunction, administrative sanction, exclusion, rescission, or internal recourse. \\
\bottomrule
\end{tabularx}
\end{table}

Legislation has moved fastest on the last two questions. The European Union has prohibited specified practices, imposed obligations proportionate to enumerated risk categories, and distributed compliance roles across providers, deployers, importers, distributors, and product manufacturers.\lrfootnote{See \statcite{euAIAct2024}.} China's Interim Measures govern publicly offered generative-AI services through provider duties concerning content, data, user protection, and security.\lrfootnote{See \statcite{chinaGenAIMeasures2023}.} In the United States, federal agencies and states have regulated particular activities, populations, and sectors, while the NIST AI Risk Management Framework supplies a voluntary vocabulary for mapping, measuring, and managing risk.\lrfootnote{See \bluecite{nistAIRMF2023}. For a proposal to understand AI law as the law of risky agents without intentions, see \artcite{ayresBalkin2024}.}

This activity is substantial. It nevertheless leaves the first two questions fragmented. A statute can allocate ``provider'' obligations without deciding whether generated output is a manifestation in contract. A consumer regulator can demand a chatbot disclosure without deciding whether the system is a product for strict-liability purposes. A copyright office can require human authorship without deciding whether a conversational interface conveys protected human speech. A procurement rule can demand security controls without deciding whether the government may retaliate against a supplier for refusing another use. The fact that each institution has sensibly defined its own trigger does not provide the common premise that adjudication needs.

AI law is thus best understood as a coordination field. Its unity lies in recurrent relations among people, organizations, automated instruments, and institutions, not in a single substantive command. Scholarship has accordingly focused on risk distribution, delegation, accidents, opacity, and the displacement of intent.\lrfootnote{See \artcite{abrahamSharkey2027}; \artcite{duffourcGerke2023}; \artcite{chanAIGovernance2021}; \bluecite{ramakrishnanRand2024}.} Those accounts demonstrate why ordinary doctrine requires translation. The translation must begin by distinguishing the regulated system from the legal subjects whose relations a court can change.

That task is especially important in federal statutory interpretation. \emph{Loper Bright} requires courts to exercise independent judgment about a statute's best reading; ambiguity alone no longer activates \emph{Chevron}'s command to defer to a permissible agency interpretation.\lrfootnote{\casecite[394--96]{loperBright2024}.} The decision did not make technical agencies irrelevant. Their expertise may persuade under \emph{Skidmore}, and a statute may expressly confer discretion within judicially policed boundaries.\lrfootnote{See \casecite[140]{skidmore1944}; \casecite[395]{loperBright2024}.} It does mean that a court cannot treat an administrative definition, designed for an agency's program, as a complete jurisprudential answer. The court must identify the function the definition serves and interpret the operative legal relation for itself.

\subsection{The Strong Case for Definition Pluralism---and Its Limit}

The strongest objection to a jurisprudential minimum is not that definitions are unimportant. It is that a single definition is likely to be both brittle and distorting. Modern systems combine statistical models, conventional software, retrieval, tools, human workflows, and changing post-deployment configurations. A definition keyed to a fashionable architecture will expire. A definition keyed to breadth may capture spreadsheets, ordinary search, or any rule-based program. A definition keyed to high capability can generate perverse incentives around benchmark boundaries. Most importantly, the conduct relevant to a medical-device regulator is not the conduct relevant to a copyright court.

Risk-based and context-specific definitions address these concerns. Legal definitions clarify propositions within an institutional practice rather than discover the abstract essence of an object.\lrfootnote{See H.L.A. Hart, \artcite{hartDefinition1954}.} Bryan Casey and Mark Lemley show how a general robot definition can become overinclusive, underinclusive, and obsolete as technical features migrate into ordinary objects.\lrfootnote{See \artcite{caseyLemleyRobot2020}.} Jonas Schuett, in turn, shows how regulatory definitions can be designed around risks and institutional purposes.\lrfootnote{See \artcite{schuettDefinition2023}.} Technology-neutral drafting likewise seeks durable rules by describing functions and effects while avoiding dependence on a transient implementation.\lrfootnote{See \artcite{ojanenTechNeutrality2025}.} That technique cannot decide which functions or effects law should regulate; it preserves a policy choice rather than eliminating one. Recent definitional work catalogs the considerable variation among enacted and proposed AI laws and warns that recycled language can carry consequences across programs for which it was never designed.\lrfootnote{See \bluecite{aspenDefinitions2026}.} These are reasons to preserve doctrinal pluralism, not to eliminate it.

But pluralism has a boundary. Three different kinds of legal statements are often treated as though all were ``AI definitions'':

\begin{enumerate}[leftmargin=2.2em]
\item A \term{regulatory trigger} specifies which systems, actors, or uses fall within a statute or administrative program.
\item A \term{functional classification} connects facts about a system to an element of a doctrine, such as product, service, instrument, publication, or decision procedure.
\item A \term{subject-status rule} identifies who can occupy a legal position in the relation---who holds the claim, owes the duty, exercises the power, bears the liability, and answers the remedy.
\end{enumerate}

The first two should vary with purpose. The third cannot drift silently with them. If an AI Act narrows its definition to exclude a low-risk optimization system, that exclusion determines statutory coverage; it does not turn the system into a legal person for tort or contract. If product law includes a conversational application because its mass-distributed design is alleged to be defective, that classification does not grant the application assets or procedural capacity. If a constitutional court recognizes human expression transmitted with generative assistance, it protects the relevant speaker without personifying the medium.

The proposal here is therefore not a universal taxonomy of AI. It is a relational invariant. Unless positive law creates and equips a new juridical person, the system is not an independent legal subject. It may be many kinds of legal object, and the people or organizations around it may occupy several regulated roles. This modest rule leaves sector-specific definitions free to evolve while preventing those definitions from becoming accidental grants of personhood or accidental releases of responsibility.

Two tests follow. First, a court should ask whether a proposed label tracks an element of the operative law. ``Product'' matters when a rule turns on distribution, design, warning, or sale; ``speaker'' matters when the government regulates expression; ``agent'' matters in agency law only if the elements of assent, acting on behalf, and control are satisfied. Second, the court should ask whether the label is being used to relocate a legal position. A term that begins as a description of functionality cannot, without an attribution rule, make the machine the promisor, tortfeasor, author, fiduciary, or constitutional rights holder.

\subsection{Private Definitions Are Evidence, Not Constitutive Law}

Private firms are unusually powerful definers of AI because they design both the system and the user's encounter with it. Interface copy assigns a name, voice, biography, and relational role. Terms of service characterize output as information, entertainment, professional assistance, or user content. System specifications describe desired behavior. Model constitutions state internal principles. Warnings characterize limitations. Together these materials may shape use before any public institution has classified the service.

They are legally relevant for that reason. Marketing can establish an intended use. A named and emotionally responsive ``companion'' can make relational reliance foreseeable in a way that a bare text-completion endpoint does not. A claim that a system can act autonomously for a customer may shape reasonable expectations about monitoring and authorization. Warnings may affect reliance, breach, or an available defense. Internal evaluations and deployment criteria may show what a provider knew, when it knew it, and which precautions were feasible.

But none of these acts is privately constitutive of legal personality. A model constitution is an organizational control document, not a constitution in the public-law sense. A term granting the user rights in output can allocate whatever rights a party possesses; it cannot create federal copyright where the Copyright Act supplies none. A disclaimer can affect a claim only through the law governing assent, public policy, causation, warranty, or duty; the phrase ``AI-generated'' is not a freestanding immunity. And calling software an ``agent'' cannot transfer liability to an entity that has no separate assets, service of process, or legally recognized obligation.

The British Columbia Civil Resolution Tribunal's decision in \emph{Moffatt v. Air Canada} provides a compact illustration. Air Canada argued, in substance, that information from its website chatbot should not be attributed to the airline. The tribunal refused to treat the chatbot as a separate responsible entity: the website was the airline's channel, and the airline was responsible for information delivered through it.\lrfootnote{\casecite[27--31]{moffatt2024}. The tribunal resolved a consumer dispute under British Columbia law; the decision is comparative evidence of an attribution problem, not controlling United States authority.} The point is not that every generated sentence invariably binds a provider. The Hangzhou court reached a different contract result because the bot's extravagant offer was not the provider's manifestation. The common principle is prior to the merits: the system does not become a third person who absorbs the legal question. The court asks whether and how the output is attributable to an existing subject under the governing rule.

This distinction also aligns incentives. If branding were constitutive, a provider could advertise autonomy when autonomy sells and insist on mere-tool status when responsibility attaches. If disclaimers could create a machine intermediary, firms could externalize risk through interface design. Treating private definitions as evidence avoids both results. It gives marketing and technical documentation their proper factual force while reserving legal status to public law.

\subsection{One System, Several Functional Classifications}

Courts already assign the same object different classifications across legal regimes. A mobile phone may be a product in tort, a service platform in contract, a protected repository in search law, and a medium of expression under the First Amendment. AI introduces harder facts, but no new requirement that one noun govern every dispute. The relevant classifications can be organized as follows:

\begin{table}[htbp]
\centering
\caption{Functional Classification Without Artificial Personhood}
\label{tab:functional-classification}
\small
\begin{tabularx}{\textwidth}{>{\raggedright\arraybackslash}p{0.17\textwidth} >{\raggedright\arraybackslash}p{0.27\textwidth} X}
\toprule
Legal setting & Classification question & What remains invariant \\
\midrule
Products and services & Is the alleged harm connected to mass-distributed design, a component, a service performance, or professional judgment? & The provider, deployer, seller, professional, and claimant remain the possible legal subjects. \\
Contract & Is output an attributable manifestation, an automated transaction, information, or conduct outside configured authority? & The system does not itself acquire contracting capacity merely by generating assent-like text. \\
Copyright & Which expressive elements reflect human authorship, and which are generated or unprotectable? & AI can be an instrument; ownership begins only with rights that law recognizes in an author or other rights holder. \\
Speech & Whose expression is regulated, and does the system transmit, transform, select, or originate the contested material? & Protection, if any, belongs to a recognized speaker; calling output ``speech'' need not make the model a rights holder. \\
\bottomrule
\end{tabularx}
\end{table}
\clearpage

Product classification demonstrates the method. The Restatement's product rules attend to commercial distribution, design, warning, and physical or other legally cognizable harm.\lrfootnote{See \bluecite{restatementProducts1998}.} Software and mixed services complicate the boundary, and courts need not resolve that boundary for every AI system at once. In \emph{Garcia}, the court held only that the complaint plausibly alleged a product at the pleading stage, relying on allegations that the application was standardized, distributed to consumers, and designed to generate the interactions at issue.\lrfootnote{\casecite{garcia2025}.} That determination selected a body of law for testing alleged enterprise conduct. It did not turn the bot into either a human speaker or an independent tortfeasor.

Contract law supplies an even older model of automatic instrumentality. The Uniform Electronic Transactions Act, a model act enacted in most states, defines an ``electronic agent'' as an automated means used to initiate an action or respond to records without contemporaneous review by an individual. It nevertheless attributes the resulting transaction through rules governing the parties; it does not confer contractual personality on the program. Federal law makes the attribution point explicit: a transaction is not denied legal effect solely because electronic agents participated, so long as the agent's action is legally attributable to the person to be bound.\lrfootnote{\bluecite{ueta1999} (model act \S\S~2(6), 14); \statcite{esignElectronicAgents}.} Generative systems make attribution more fact intensive because outputs are open ended and authority may be constrained through prompts, tools, spending limits, confirmations, and interface design. The legal question remains whether an existing party manifested assent or authorized an automated transaction.

Copyright likewise separates instrument from author. In \emph{Thaler v. Perlmutter}, the D.C. Circuit held that the Copyright Act requires a human author, while emphasizing that human use of AI does not categorically defeat protection for human-authored contributions.\lrfootnote{\casecite[1045--49]{thaler2025}; see also \bluecite{uscoPartTwo2025}.} The Copyright Office's analysis asks which expressive choices are attributable to a human, not whether the system is sophisticated enough to deserve a share. This rule coexists comfortably with a product classification in another case.

Speech questions require the same care. Generated text can contain a user's expression, embody a provider's editorial or product-design choices, reproduce third-party material, or emerge without any human adopting the particular words.\lrfootnote{See \artcite{jonesKaminskiAISpeech2024}.} The First Amendment analysis must identify the claimant, the regulated conduct, and the expressive interest. The \emph{Garcia} court therefore declined to announce that Character.AI output was protected speech ``at this stage'' when the defendants had not identified a human speaker whose rights would be burdened.\lrfootnote{\casecite{garcia2025}.} That ruling leaves room for a different record in which a user's AI-assisted composition or a firm's editorial design carries a recognizable human claim.

The result is disciplined pluralism. Courts may select the classification that makes the operative doctrine work, explain the factual feature that justifies it, and confine the conclusion to that function. What they should not do is let a functional label alter the legal cast of characters without positive law. Part II now develops that cast more precisely.
